\documentclass[17pt,compress]{beamer}
\usepackage{beamerthemesplit}
\mode<presentation>
{
  \usetheme{Warsaw}
  \setbeamercovered{transparent}
  \setbeamertemplate{navigation symbols}{}
}
% Taken from Fernando's slides.
\usepackage{ae,aecompl}
\usepackage[scaled=.95]{helvet}

\usepackage[english]{babel}
\usepackage[latin1]{inputenc}
\usepackage[T1]{fontenc}

\newcounter{saveenumi}
\newcommand{\seti}{\setcounter{saveenumi}{\value{enumi}}}
\newcommand{\conti}{\setcounter{enumi}{\value{saveenumi}}}

\definecolor{blue}{rgb}{0.16,0.32,0.75}
\setbeamercolor{structure}{fg=blue}
\author[FOSSEE]{}
\institute[IIT Bombay]{}
\date[]{}
% \setbeamercovered{transparent}

% theme split
\usepackage{verbatim}
\newenvironment{colorverbatim}[1][]%
{%
\color{blue}
\verbatim
}%
{%
\endverbatim
}%

\usepackage{mathpazo,courier,euler}
\usepackage{listings}
\lstset{language=sh,
    basicstyle=\ttfamily\bfseries,
  showstringspaces=false,
  keywordstyle=\color{black}\bfseries}

% logo
\logo{\includegraphics[height=1.30 cm]{3t-logo.pdf}}
\logo{\includegraphics[height=1.30 cm]{fossee-logo.png}

\hspace{7.5cm}
\includegraphics[scale=0.3]{fossee-logo.png}\\
\hspace{281pt}
\includegraphics[scale=0.80]{3t-logo.pdf}}

\begin{document}

\sffamily \bfseries
\title
[Using Python Modules]
{Python - Using Python Modules}
\author
[FOSSEE, IIT - Bombay]
{\small Talk to a Teacher\\{\color{blue}\url{http://spoken-tutorial.org}}\\National Mission on Education
 through ICT\\{\color{blue}\url{http://sakshat.ac.in}} \\[0.5cm]{\tiny Script by: Hardik ghaghada \\ Naration by: Hardik Ghaghada \\ 13 August 2015}}

% slide 1
\begin{frame}
   \titlepage
\end{frame}
%%%%%%%%%%%%%%%%%%%%%%%%%%%%%%%%%%%%%%%%%%%%%%%%%%%%%%%%%%%%%%%%%%%%%%%%%%%%%%%%
\begin{frame}
\frametitle{Objectives}
\label{sec-2}

  In this tutorial we will learn, \pause

\begin{itemize}
\item Execute python scripts from command line.\pause
\item Use import in scripts.\pause
\item Import scipy and pylab modules.\pause
\item Use python standard modules and 3rd party modules.
\end{itemize}
\end{frame}
%%%%%%%%%%%%%%%%%%%%%%%%%%%%%%%%%%%%%%%%%%%%%%%%%%%%%%%%%%%%%%%%%%%%%%%%%%%%%%%%
\begin{frame}
\frametitle{System Specifications}\pause
\begin{itemize}
\item Slackware Linux 14.1\pause
\item \texttt{Python 2.7.5} \pause
\item \texttt{IPython 2.3.0}
\end{itemize}
\end{frame}
%%%%%%%%%%%%%%%%%%%%%%%%%%%%%%%%%%%%%%%%%%%%%%%%%%%%%%%%%%%%%%%%%%%%%%%%%%%%%%%%
\begin{frame}
\frametitle{Pre-requisites}
\label{sec-3}

Spoken tutorial on -

\begin{itemize}
\item Using plot interactively.\pause
\item Embellishing a plot.\pause
\item Saving plots.
\end{itemize}
\end{frame}
%%%%%%%%%%%%%%%%%%%%%%%%%%%%%%%%%%%%%%%%%%%%%%%%%%%%%%%%%%%%%%%%%%%%%%%%%%%%%%%%
\begin{frame}[fragile]
\frametitle{Running Python script from command line}
\label{sec-4}

\begin{itemize}
\item Create a script, open text editor and type the following\\\pause
     \texttt{print "hello world!"}\\
     \texttt{print}\pause
\item Save the script as \texttt{hello.py}
\end{itemize}
\end{frame}
%%%%%%%%%%%%%%%%%%%%%%%%%%%%%%%%%%%%%%%%%%%%%%%%%%%%%%%%%%%%%%%%%%%%%%%%%%%%%%%%
\begin{frame}[fragile]
\frametitle{Running Python script from command line (cont'd)}
\label{sec-5}


\begin{itemize}
\item Run the script\\
     \texttt{\$ python hello.py}
\end{itemize}
  \emph{Syntax :} \textbf{python filename}
\end{frame}
%%%%%%%%%%%%%%%%%%%%%%%%%%%%%%%%%%%%%%%%%%%%%%%%%%%%%%%%%%%%%%%%%%%%%%%%%%%%%%%%
\begin{frame}
\frametitle{Four plot problem}
\label{sec-6}

    \begin{center}
      \includegraphics[scale=0.4]{four_plot}    
    \end{center}
\end{frame}
%%%%%%%%%%%%%%%%%%%%%%%%%%%%%%%%%%%%%%%%%%%%%%%%%%%%%%%%%%%%%%%%%%%%%%%%%%%%%%%%
\begin{frame}[fragile]
\frametitle{Fix \texttt{linspace()} problem}
\label{sec-7}

\texttt{from scipy import *}
\end{frame}
%%%%%%%%%%%%%%%%%%%%%%%%%%%%%%%%%%%%%%%%%%%%%%%%%%%%%%%%%%%%%%%%%%%%%%%%%%%%%%%%
\begin{frame}[fragile]
\frametitle{Fix \texttt{plot()} problem}
\label{sec-8}
\texttt{from pylab import *}
\end{frame}
%%%%%%%%%%%%%%%%%%%%%%%%%%%%%%%%%%%%%%%%%%%%%%%%%%%%%%%%%%%%%%%%%%%%%%%%%%%%%%%%
\begin{frame}[fragile]
\frametitle{Better way of fixing}
\label{sec-9}

\texttt{from scipy import linspace}\\\pause
\vspace{5pt}
  instead of\\\pause
\vspace{5pt}
\texttt{from scipy import *}\\\pause
\vspace{5pt}
\texttt{*} means import all functions from name-space \texttt{scipy}.
\end{frame}
%%%%%%%%%%%%%%%%%%%%%%%%%%%%%%%%%%%%%%%%%%%%%%%%%%%%%%%%%%%%%%%%%%%%%%%%%%%%%%%%
\begin{frame}[fragile]
\frametitle{Instead of \texttt{*}}
\label{sec-10}
\begin{small}
\texttt{from scipy import linspace, pi, sin}
\texttt{from pylab import plot, legend, annotate}
  
\texttt{from pylab import xlim, ylim, title, show}\pause
\end{small}\\
  Is better than,\\ {\small \texttt{from scipy import *} \& \texttt{from pylab import *}}.
\end{frame}
%%%%%%%%%%%%%%%%%%%%%%%%%%%%%%%%%%%%%%%%%%%%%%%%%%%%%%%%%%%%%%%%%%%%%%%%%%%%%%%%
\begin{frame}[fragile]
\frametitle{Another Fix}
\label{sec-11}

\lstset{language=Python}
\begin{footnotesize}
\begin{lstlisting}
import scipy
import pylab
x = scipy.linspace(-5*scipy.pi,5*scipy.pi,500)
pylab.plot(x, x, 'b')
pylab.plot(x, -x, 'b')
pylab.plot(x, scipy.sin(x), 'g', linewidth=2)
pylab.plot(x, x*scipy.sin(x), 'r', linewidth=3)
pylab.legend(['x', '-x', 'sin(x)', 'xsin(x)'])
pylab.annotate('origin', xy = (0, 0))
pylab.xlim(-5*scipy.pi, 5*scipy.pi)
pylab.ylim(-5*scipy.pi, 5*scipy.pi)
\end{lstlisting}
\end{footnotesize}
\end{frame}
%%%%%%%%%%%%%%%%%%%%%%%%%%%%%%%%%%%%%%%%%%%%%%%%%%%%%%%%%%%%%%%%%%%%%%%%%%%%%%%%
\begin{frame}
\frametitle{Exercise 1}
\label{sec-12}

\begin{itemize}
\item Write a python script to plot a sine wave from 
    $-2\pi$
  to 
    $2\pi$
  .
\end{itemize}
\end{frame}
%%%%%%%%%%%%%%%%%%%%%%%%%%%%%%%%%%%%%%%%%%%%%%%%%%%%%%%%%%%%%%%%%%%%%%%%%%%%%%%%
\begin{frame}
\frametitle{What is a module?}
\label{sec-13}
\begin{itemize}
\item Module is simply a file containing Python definitions and
  statements.\pause
\item Definitions from a module can be imported into other
  modules or into the main module.
\end{itemize}
  
\end{frame}
%%%%%%%%%%%%%%%%%%%%%%%%%%%%%%%%%%%%%%%%%%%%%%%%%%%%%%%%%%%%%%%%%%%%%%%%%%%%%%%%
\begin{frame}
\frametitle{Python standard library}
\label{sec-14.1}

  Python has a very rich standard library of modules.\pause

\begin{itemize}
\item Few libraries
	\begin{itemize}
	\item Math: \texttt{math}, \texttt{random}\pause
	\item Internet access: \texttt{urllib2}, \texttt{smtplib}\pause
	\item System, Command line arguments: \texttt{sys}
	\end{itemize}
\end{itemize}
\end{frame}
%%%%%%%%%%%%%%%%%%%%%%%%%%%%%%%%%%%%%%%%%%%%%%%%%%%%%%%%%%%%%%%%%%%%%%%%%%%%%%%%
\begin{frame}
\frametitle{Python standard library}
\label{sec-14.2}

\begin{itemize}
\item Few more libraries
	\begin{itemize}
	\item Operating system interface: \texttt{os}\pause
	\item regular expressions: \texttt{re}\pause
	\item compression: \texttt{gzip}, \texttt{zipfile}, \texttt{tarfile}
	\end{itemize}\pause
\item More information
	\begin{itemize}
	\item \href{http://docs.python.org/library}{http://docs.python.org/library}
	\end{itemize}
\end{itemize}
\end{frame}
%%%%%%%%%%%%%%%%%%%%%%%%%%%%%%%%%%%%%%%%%%%%%%%%%%%%%%%%%%%%%%%%%%%%%%%%%%%%%%%%
\begin{frame}
\frametitle{Summary}
\label{sec-15.1}

  In this tutorial, we have learnt to -

\begin{itemize}
\item Run scripts from command line,
\item Import modules by specifying the module name followed by  
    an asterisk.
\end{itemize}
\end{frame}
%%%%%%%%%%%%%%%%%%%%%%%%%%%%%%%%%%%%%%%%%%%%%%%%%%%%%%%%%%%%%%%%%%%%%%%%%%%%%%%%
\begin{frame}
\frametitle{Summary contd..}
\label{sec-15.2}

\begin{itemize}
\item Import only the required functions from modules by specifying 
    the function name.
\item Use python standard library.
\end{itemize}
\end{frame}
%%%%%%%%%%%%%%%%%%%%%%%%%%%%%%%%%%%%%%%%%%%%%%%%%%%%%%%%%%%%%%%%%%%%%%%%%%%%%%%%
\begin{frame}
\frametitle{Evaluation}
\label{sec-16.1}

\begin{enumerate}
\item Which among this is correct ?\pause
	\begin{itemize}
	\item \texttt{from scipy import plot}
	\item \texttt{from numpy import plot}
	\item \texttt{from matplotlib import plot}
	\item \texttt{from pylab import plot}
	\end{itemize}\pause
\seti
\end{enumerate}
\end{frame}
%%%%%%%%%%%%%%%%%%%%%%%%%%%%%%%%%%%%%%%%%%%%%%%%%%%%%%%%%%%%%%%%%%%%%%%%%%%%%%%%
\begin{frame}
\frametitle{Evaluation contd..}
\label{sec-16.2}

\begin{enumerate}
\item Which among these libraries is part of python standard library ?\pause
	\begin{itemize}
	\item \texttt{Mayavi}
	\item \texttt{scipy}
	\item \texttt{matplotlib}
	\item \texttt{urllib2}
	\end{itemize}
\seti
\end{enumerate}
\end{frame}
%%%%%%%%%%%%%%%%%%%%%%%%%%%%%%%%%%%%%%%%%%%%%%%%%%%%%%%%%%%%%%%%%%%%%%%%%%%%%%%%
\begin{frame}
\frametitle{Evaluation contd..}
\label{sec-16.3}

\begin{enumerate}
\conti
\item Functions \texttt{xlim()} and \texttt{ylim()} can be imported to the current
   name-space as,\pause
	\begin{itemize}
	\item \texttt{from pylab import xlim, ylim}
	\item \texttt{import pylab}
	\item \texttt{from scipy import xlim, ylim}
	\item \texttt{import scipy}
	\end{itemize}
\end{enumerate}
\end{frame}
%%%%%%%%%%%%%%%%%%%%%%%%%%%%%%%%%%%%%%%%%%%%%%%%%%%%%%%%%%%%%%%%%%%%%%%%%%%%%%%%
\begin{frame}
\frametitle{Solutions}
\label{sec-17}

\begin{enumerate}
\item \texttt{from pylab import plot}\pause
\vspace{12pt}
\item \texttt{urllib2}\pause
\vspace{12pt}
\item \texttt{from pylab import xlim, ylim}
\end{enumerate}
\end{frame}
%%%%%%%%%%%%%%%%%%%%%%%%%%%%%%%%%%%%%%%%%%%%%%%%%%%%%%%%%%%%%%%%%%%%%%%%%%%%%%%%
\begin{frame}
\frametitle{FOSSEE}
{\color{blue}Free and Open-source Software for \\Science and Engineering Education} \\
\begin{itemize}
\item Goal - enabling all to use open source software tools
\item For more details, please visit {\color{blue}\url{http://fossee.in/}}
\end{itemize}
\end{frame}
%%%%%%%%%%%%%%%%%%%%%%%%%%%%%%%%%%%%%%%%%%%%%%%%%%%%%%%%%%%%%%%%%%%%%%%%%%%%%%%%
\begin{frame}
\frametitle{About the Spoken Tutorial Project}
\begin{itemize}
\item Watch the video available at {\color{blue}http://spoken-tutorial.org /What\_is\_a\_Spoken\_Tutorial}
\item It summarises the Spoken Tutorial project \pause
\item If you do not have good bandwidth, you can download and watch it
\end{itemize}
\end{frame}
%%%%%%%%%%%%%%%%%%%%%%%%%%%%%%%%%%%%%%%%%%%%%%%%%%%%%%%%%%%%%%%%%%%%%%%%%%%%%%%%
\begin{frame}
\frametitle{Spoken Tutorial Workshops}The Spoken Tutorial Project Team 
\begin{itemize}
\item Conducts workshops using spoken tutorials 
\item Gives certificates to those who pass an online test 
\item For more details, please write to \\ \hspace {0.5cm}{\color{blue}contact@spoken-tutorial.org}
\end{itemize}
\end{frame}
%%%%%%%%%%%%%%%%%%%%%%%%%%%%%%%%%%%%%%%%%%%%%%%%%%%%%%%%%%%%%%%%%%%%%%%%%%%%%%%%
\begin{frame}
\frametitle{Acknowledgements}
\begin{itemize}
\item Spoken Tutorial Project is a part of the Talk to a Teacher  project 
\item It is supported by the National Mission on Education through  ICT, MHRD, Government of India 
\item More information on this Mission is available at: \\{\color{blue}http://spoken-tutorial.org /NMEICT-Intro}
\end{itemize}
\end{frame}
%%%%%%%%%%%%%%%%%%%%%%%%%%%%%%%%%%%%%%%%%%%%%%%%%%%%%%%%%%%%%%%%%%%%%%%%%%%%%%%%
\begin{frame}

  \begin{block}{}
  \begin{center}
  \textcolor{blue}{\Large THANK YOU!} 
  \end{center}
  \end{block}
\begin{block}{}
  \begin{center}
    For more Information, visit our website\\
    \url{http://fossee.in/}
  \end{center}  
  \end{block}
\end{frame}
%%%%%%%%%%%%%%%%%%%%%%%%%%%%%%%%%%%%%%%%%%%%%%%%%%%%%%%%%%%%%%%%%%%%%%%%%%%%%%%%
\end{document}
